% © 2025 Ing. David Jaroš — CC BY-NC-ND 4.0
%
% This work is licensed under a Creative Commons Attribution-NonCommercial-NoDerivatives
% 4.0 International License (CC BY-NC-ND 4.0).
%
% License History: Earlier drafts (up to v0.3) were released under CC BY 4.0.
% From v0.4 onward, all material is released under CC BY-NC-ND 4.0 to protect
% the integrity of the theoretical work during ongoing academic development.

% =====================================================================
% GAP 3 — Compute ∂α/∂φ: Is φ Physical or Gauge?
% Target: canonical/geometry/phi_gauge_vs_physical.tex
% Status: [DERIVED] — conditional on UBT vacuum choice (see §4)
% Layer: [L1] — biquaternionic geometry
% =====================================================================

% UBT-AI-PROVENANCE-BEGIN
% schema: ubt-ai-provenance/v1
% tier: B_machine_verified
% ai_assistance: disclosed
% human_review: machine-verification
% editorial_responsibility: Ing. David Jaroš
% policy: ../../AI_PROVENANCE.md
% notice: Machine-verified against named sources or verifiers; individual attestation is not claimed.
% UBT-AI-PROVENANCE-END

\section{Is $\varphi$ Physical or Pure Gauge? — Computation of $\partial\alpha/\partial\varphi$}
\label{sec:phi_gauge_vs_physical}

\subsection{Overview}

The biquaternionic metric $\mathcal{U}_{\mu\nu}$ admits a one-parameter family of real
projections indexed by a phase angle $\varphi \in [0, 2\pi)$:
\begin{equation}
  P_\varphi[\mathcal{U}_{\mu\nu}] = \mathrm{Re}(e^{-i\varphi} \cdot \mathcal{U}_{\mu\nu}).
  \label{eq:phase_projection}
\end{equation}
\textbf{[POSTULATE — see \texttt{canonical/geometry/biquaternion\_metric.tex}]}

\noindent \textbf{Key question:} Is $\varphi$ a gauge redundancy (all frames physically identical)
or a genuine moduli parameter (different $\varphi$ $\Rightarrow$ different physics)?

\noindent \textbf{Diagnostic:} Compute $\partial\alpha/\partial\varphi$.  If it is zero,
$\varphi$ is pure gauge.  If nonzero, $\varphi$ is physical.

\noindent \textbf{Note on notation:} Throughout, $\varphi$ denotes the phase-frame parameter in
$P_\varphi[\mathcal{U}_{\mu\nu}]$ (observer choice), and $\psi$ is the imaginary part of complex
time $\tau = t + i\psi$ (dynamical field).  These are distinct objects and must not be conflated.

---

\subsection{Step 1 — Phase Rotation of $\mathcal{U}_{\mu\nu}$}
\label{sec:phi_step1}

\textbf{[POSTULATE]} From \texttt{canonical/geometry/biquaternion\_metric.tex},
the biquaternionic metric decomposes as:
\begin{equation}
  \mathcal{U}_{\mu\nu} = g_{\mu\nu} + i\,h_{\mu\nu} + j\,k^1_{\mu\nu} + k\,k^2_{\mu\nu},
  \label{eq:biq_metric_decomp}
\end{equation}
where $g_{\mu\nu}, h_{\mu\nu}, k^1_{\mu\nu}, k^2_{\mu\nu} \in \mathbb{R}$ are the real,
imaginary-scalar, and quaternionic components respectively.

Applying the U(1) phase rotation $e^{-i\varphi}$ (scalar complex rotation,
leaving the quaternionic basis $\{j, k\}$ unchanged):
\begin{align}
  e^{-i\varphi} \cdot \mathcal{U}_{\mu\nu}
  &= e^{-i\varphi}(g_{\mu\nu} + i\,h_{\mu\nu}) + e^{-i\varphi}(j\,k^1_{\mu\nu} + k\,k^2_{\mu\nu})
  \notag\\
  &= \bigl(\cos\varphi\,g_{\mu\nu} + \sin\varphi\,h_{\mu\nu}\bigr)
   + i\bigl(-\sin\varphi\,g_{\mu\nu} + \cos\varphi\,h_{\mu\nu}\bigr)
   + j\,k^1_{\mu\nu}\cos\varphi + \ldots
\end{align}
Taking the real part:
\begin{equation}
  \boxed{P_\varphi[\mathcal{U}_{\mu\nu}]
    = \cos\varphi \cdot g_{\mu\nu} + \sin\varphi \cdot h_{\mu\nu}.}
  \label{eq:proj_result}
  \quad \textbf{[DERIVED]}
\end{equation}
Note: the quaternionic components $\{k^1, k^2\}$ are unaffected by the scalar U(1) rotation.

---

\subsection{Step 2 — EM Gauge Potential Under $\varphi$-Rotation}
\label{sec:phi_step2}

U(1)$_\mathrm{EM}$ comes from the complex phase of $\Theta$ (proved in
\texttt{canonical/interactions/qed.tex}).  Write the biquaternionic gauge potential as:
\begin{equation}
  \mathcal{A}_\mu = \mathcal{A}^R_\mu + i\,\mathcal{A}^I_\mu, \quad
  \mathcal{A}^R_\mu, \mathcal{A}^I_\mu \in \mathbb{R}^4.
\end{equation}
After $\varphi$-projection (same U(1) rotation):
\begin{equation}
  A^{(\varphi)}_\mu = \cos\varphi \cdot \mathcal{A}^R_\mu + \sin\varphi \cdot \mathcal{A}^I_\mu.
  \label{eq:A_phi}
  \quad \textbf{[DERIVED]}
\end{equation}

---

\subsection{Step 3 — $\alpha(\varphi)$ Formula}
\label{sec:phi_step3}

The coupling strength $e(\varphi)^2 \propto \langle A^{(\varphi)}, A^{(\varphi)}\rangle$.
Define the amplitude ratio and correlation coefficient:
\begin{align}
  r &:= \frac{|\mathcal{A}^I_\mu|}{|\mathcal{A}^R_\mu|}, \qquad
  \rho := \frac{\mathrm{Re}(\mathcal{A}^R \cdot \mathcal{A}^I)}{|\mathcal{A}^R||\mathcal{A}^I|}.
\end{align}
Then:
\begin{align}
  |A^{(\varphi)}|^2
  &= \cos^2\!\varphi\,|\mathcal{A}^R|^2
   + 2\rho\,r\cos\varphi\sin\varphi\,|\mathcal{A}^R|^2
   + r^2\sin^2\!\varphi\,|\mathcal{A}^R|^2 \notag\\
  &= |\mathcal{A}^R|^2\bigl[\cos^2\!\varphi + 2\rho r\cos\varphi\sin\varphi + r^2\sin^2\!\varphi\bigr].
\end{align}
Since $\alpha \propto e^2 \propto |A^{(\varphi)}|^2 / \text{(normalization)}$:
\begin{equation}
  \boxed{\alpha(\varphi) = \alpha(0)\bigl[\cos^2\!\varphi + 2\rho r\cos\varphi\sin\varphi + r^2\sin^2\!\varphi\bigr].}
  \label{eq:alpha_phi}
  \quad \textbf{[DERIVED]}
\end{equation}

Taking the derivative:
\begin{equation}
  \frac{\partial\alpha}{\partial\varphi}\bigg|_{\varphi=0}
  = \alpha(0)\bigl[2\rho r - 0\bigr]
  = 2\rho r\,\alpha(0).
  \label{eq:dalpha_dphi}
  \quad \textbf{[DERIVED]}
\end{equation}

\noindent \textbf{Therefore:}
\begin{itemize}
  \item The analytic criterion is $\rho r=0$ versus $\rho r\neq0$.
  \item A real metric component $h_{\mu\nu}=0$ does not by itself prove
    $r=0$ for an independently defined gauge potential.
  \item Conversely, $h_{\mu\nu}\neq0$ does not by itself prove either
    $r\neq0$ or $\rho\neq0$.  The map from the metric sector to the
    canonically normalized gauge potential must be derived separately.
\end{itemize}

---

\subsection{Step 4 --- Evaluation for the UBT Vacuum}
\label{sec:phi_step4}

\subsubsection{Flat branch}

For a branch in which the canonical gauge potential is purely real,
$r=0$ and hence
\begin{equation}
 \frac{\partial\alpha}{\partial\varphi}=0.
\end{equation}
This is a conditional statement about the gauge sector, not a consequence of
$h_{\mu\nu}=0$ alone.

\subsubsection{Nontrivial biquaternionic branch}

A physical phase modulus requires an explicit solution for which the
canonically derived gauge components satisfy
\begin{equation}
 r\neq0,\qquad \rho\neq0.
\end{equation}
Then and only then does
$\partial\alpha/\partial\varphi=2\rho r\alpha(0)$ establish local physical
phase dependence.  The former two-mode Hermitian example does not provide
such a solution: its metric component is real, its gauge-potential expression
is only a sketch, and the reported correlation is $\rho\simeq0$.

\noindent\textbf{Current conclusion.}
The analytic diagnostic is derived, but its nonzero realization in a
canonical UBT vacuum remains open.  No prime-selected value of $\varphi$ is
predicted by this calculation alone.

---

\subsection{Summary Table}
\label{sec:phi_summary}

\begin{center}
\begin{tabular}{p{0.32\linewidth}p{0.24\linewidth}p{0.32\linewidth}}
\hline
\textbf{Evidence} & $\partial\alpha/\partial\varphi$ & \textbf{$\varphi$ status} \\
\hline
Canonical $\rho r=0$ & $0$ & Gauge/degenerate on that branch \\
Canonical $\rho r\neq0$ & $2\rho r\alpha(0)\neq0$ & Physical modulus on that branch \\
Only $h_{\mu\nu}\neq0$ or sketch $r\neq0$ & Undetermined & Open \\
\hline
\end{tabular}
\end{center}

\noindent\textbf{Layer:} [L1] --- conditional biquaternionic geometry.

\noindent\textbf{Repository correction (31 July 2026).}
The formula for $\partial\alpha/\partial\varphi$ remains derived.  The previous
claim that a two-mode Hermitian vacuum realizes a physical $\varphi$ has been
withdrawn; the realization problem is open.

